% Actividad 4 — Tabla cronológica sobre la evolución de la teoría administrativa
\documentclass[11pt,letterpaper,oneside]{article}

\usepackage[margin=2.2cm]{geometry}
\usepackage[utf8]{inputenc}
\usepackage[T1]{fontenc}
\usepackage[spanish,es-tabla]{babel}
\usepackage[scaled=.96]{helvet}
\renewcommand{\familydefault}{\sfdefault}
\usepackage{graphicx}
\usepackage[table]{xcolor}
\usepackage{booktabs}
\usepackage{array}
\usepackage{longtable}
\usepackage{ragged2e}
\usepackage{lscape}
\usepackage{hyperref}
\usepackage[authoryear,round]{natbib}

\definecolor{itescaRedDark}{HTML}{8F1117}
\definecolor{itescaRed}{HTML}{D70F18}
\definecolor{itescaCharcoal}{HTML}{141414}
\definecolor{itescaSilver}{HTML}{D9D9D9}
\definecolor{itescaPaper}{HTML}{F7F7F7}
\definecolor{itescaInk}{HTML}{181818}
\definecolor{itescaMuted}{HTML}{626262}

\hypersetup{
  pdftitle={Actividad 4. Tabla cronológica de la evolución de la teoría administrativa},
  pdfauthor={Martin Jonathan de la Cruz},
  pdfsubject={Fundamentos de Gestión Administrativa},
  colorlinks=true,
  linkcolor=itescaRedDark,
  citecolor=itescaRedDark,
  urlcolor=itescaRedDark
}
\setlength{\parindent}{0pt}
\setlength{\parskip}{0.65em}
\setlength{\LTpre}{0.15em}
\setlength{\LTpost}{0pt}
\setlength{\LTleft}{0pt}
\setlength{\LTright}{0pt}
\renewcommand{\arraystretch}{1.23}
\linespread{1.22}\selectfont
\bibliographystyle{apalike}

\newcolumntype{L}[1]{>{\RaggedRight\arraybackslash}p{#1}}
\newcommand{\teoria}[1]{\textcolor{itescaRedDark}{\textbf{#1}}}
\newcommand{\aporte}[1]{\textbullet\ #1\par}
\newcommand{\horizontalpageson}{\pdfpageattr{/Rotate 90}}
\newcommand{\horizontalpagesoff}{\pdfpageattr{}}

\begin{document}

\begin{titlepage}
  \thispagestyle{empty}
  \color{itescaInk}
  \noindent
  \begin{minipage}[t]{0.62\textwidth}
    \vspace{0pt}
    {\small Instituto Tecnológico Superior de Cajeme\par}
    {\small Maestría en Gestión Administrativa\par}
    {\small Posgrado e Investigación\par}
  \end{minipage}%
  \begin{minipage}[t]{0.38\textwidth}
    \vspace{0pt}\raggedleft
    \includegraphics[width=5.2cm]{ITESCA/assets-itesca/web/logo-itesca-contorno.png}
  \end{minipage}

  \vfill
  \begin{center}
    {\color{itescaRedDark}\fontsize{25}{30}\selectfont\bfseries
      Tabla cronológica de la evolución de la teoría administrativa\par}
    \vspace{0.45cm}
    {\large Actividad 4 -- Fundamentos de Gestión Administrativa\par}
  \end{center}

  \vfill
  \noindent
  \begin{minipage}{\textwidth}
    \raggedleft
    \renewcommand{\arraystretch}{1.18}
    \begin{tabular}{ll}
      \textbf{Alumno:} & Martin Jonathan de la Cruz \\
      No. de control: & Por confirmar \\
      Carrera: & Maestría en Gestión Administrativa \\
      Semestre: & Agosto--noviembre de 2026 \\
      Asignatura: & Fundamentos de Gestión Administrativa \\
      Docente: & Docente de la materia por confirmar \\
      Modalidad: & En línea \\
      Fecha de entrega: & 24 de agosto de 2026 \\
      Lugar: & Cajeme, Sonora
    \end{tabular}
  \end{minipage}
\end{titlepage}

% Técnica didáctica: tabla cronológica. Criterio: teoría -> periodo -> precursores -> aportaciones.
% Fuente de instrucciones: aula virtual GGFG02, Unidad II, actividad 2.2.
\section{Introducción}

La teoría administrativa no evolucionó mediante la sustitución completa de una escuela por otra. Cada enfoque respondió a problemas de su época y desplazó la atención hacia una dimensión antes insuficientemente explicada: primero la eficiencia de la tarea y la estructura formal; después las personas, la motivación y la decisión; más tarde la interacción sistémica con el ambiente; y, finalmente, la estrategia, el conocimiento y la adaptación. Por ello, las fechas de la tabla representan hitos de formulación o difusión, no límites rígidos de vigencia. En la práctica contemporánea coexisten aportes de varias escuelas \citep{larrarteSinFecha,bright2019,sanchezBanuelos2017}.

\textbf{Objetivo:} ordenar cronológicamente las principales teorías administrativas e identificar, para cada una, su periodo, sus precursores y sus aportaciones centrales. La síntesis se apoya en el recorrido didáctico de \citet{larrarteSinFecha}, se contrasta con la exposición histórica de \citet{bright2019} y se amplía hacia perspectivas organizacionales recientes con \citet{sanchezBanuelos2017}.

Los criterios de evaluación y entrega son la secuencia cronológica, la presencia de los cuatro campos solicitados, la precisión conceptual, la legibilidad de la tabla y el uso de referencias bibliográficas en formato APA.

\clearpage
\newgeometry{left=0.85cm,right=0.85cm,top=0.75cm,bottom=0.75cm}
\pagestyle{empty}
\horizontalpageson
\begin{landscape}
\phantomsection
\addcontentsline{toc}{section}{Trayectoria de la teoría administrativa}
\scriptsize
\setlength{\tabcolsep}{4pt}
\renewcommand{\arraystretch}{1.16}
\begin{longtable}{L{3.7cm} L{2.9cm} L{4.6cm} L{12.8cm}}
\caption{\textbf{Tabla cronológica: evolución de la teoría administrativa}}\label{tab:evolucion-administrativa}\\[-0.45em]
\toprule
\rowcolor{itescaCharcoal}
\textcolor{white}{\textbf{Nombre de la teoría}} &
\textcolor{white}{\textbf{Años o periodos}} &
\textcolor{white}{\textbf{Precursores}} &
\textcolor{white}{\textbf{Aportaciones principales}} \\
\midrule
\endfirsthead

\multicolumn{4}{r}{\small\itshape Continuación de la Tabla~\ref{tab:evolucion-administrativa}}\\
\toprule
\rowcolor{itescaCharcoal}
\textcolor{white}{\textbf{Nombre de la teoría}} &
\textcolor{white}{\textbf{Años o periodos}} &
\textcolor{white}{\textbf{Precursores}} &
\textcolor{white}{\textbf{Aportaciones principales}} \\
\midrule
\endhead

\midrule
\endfoot

\bottomrule
\endlastfoot

\rowcolor{itescaSilver!28}
\teoria{Administración científica} &
Finales del siglo XIX y primeras décadas del XX &
Frederick W. Taylor; Henry L. Gantt; Frank y Lillian Gilbreth; Carl G. Barth &
\aporte{Privilegió la observación, la medición y la estandarización de los métodos de trabajo.}
\aporte{Analizó las tareas mediante tiempos y movimientos para orientar su organización.}
\aporte{Concentró la atención administrativa en la tarea y en la eficiencia operacional.}
Taylor articuló el sistema; Gantt aportó programación y control, los Gilbreth desarrollaron estudios de movimientos y Barth contribuyó al cálculo y al estudio de la fatiga. El periodo indica su formación y difusión, no una fecha única de inicio \citep[sección 3.4]{bright2019}. \\

\teoria{Teoría burocrática} &
Década de 1890 y primeras décadas del siglo XX; desarrollo retomado después de 1904 &
Max Weber &
\aporte{Explicó la autoridad racional--legal como base de la organización moderna.}
\aporte{Formalizó reglas escritas, jerarquías, división técnica del trabajo, impersonalidad, mérito y continuidad de los cargos.}
\aporte{Permitió previsibilidad y trato uniforme, aunque su aplicación extrema puede producir rigidez, papeleo y lentitud.}
El aporte central fue reemplazar el favor personal por competencias y normas conocidas. El periodo corresponde al desarrollo intelectual de Weber, contemporáneo de Taylor y Fayol, no a una vigencia exclusiva de la burocracia \citep[sección 3.5]{bright2019}. \\

\rowcolor{itescaSilver!28}
\teoria{Teoría clásica o administrativa} &
1916; predominio aproximado entre 1916 y 1940 &
Henri Fayol; James D. Mooney; Lyndall Urwick; Luther Gulick &
\aporte{Trasladó el análisis desde la tarea individual hacia la organización completa y su estructura.}
\aporte{Identificó funciones empresariales y sistematizó las funciones administrativas: prever o planear, organizar, dirigir, coordinar y controlar.}
\aporte{Formuló principios como división del trabajo, autoridad y responsabilidad, unidad de mando, jerarquía, equidad y espíritu de equipo.}
Su legado permanece en el proceso administrativo y el diseño formal de responsabilidades \citep{larrarteSinFecha,bright2019}. \\

\teoria{Relaciones humanas} &
1924--1932; expansión durante las décadas de 1930 y 1940 &
Elton Mayo; Mary Parker Follett; Fritz Roethlisberger; Kurt Lewin &
\aporte{Mostró que la productividad no depende únicamente de condiciones físicas o incentivos económicos.}
\aporte{Incorporó grupos informales, comunicación, liderazgo, participación, actitudes y necesidades sociales.}
\aporte{Reubicó al trabajador como persona y miembro de grupos, cuestionando la visión exclusivamente mecánica de la organización.}
Los estudios de Hawthorne impulsaron este giro, aunque su interpretación debe leerse críticamente \citep{larrarteSinFecha,bright2019}. \\

\rowcolor{itescaSilver!28}
\teoria{Enfoque conductual o ciencias del comportamiento} &
Décadas de 1940--1960; hito didáctico en 1957 &
Herbert A. Simon; Abraham Maslow; Frederick Herzberg; Douglas McGregor; Chester Barnard &
\aporte{Profundizó el estudio de motivación, liderazgo, aprendizaje, conducta individual y dinámica organizacional.}
\aporte{Simon situó la toma de decisiones y la racionalidad limitada en el centro de la administración.}
\aporte{Maslow, Herzberg y McGregor ofrecieron marcos para comprender necesidades, satisfacción y supuestos directivos sobre las personas.}
La administración pasó de prescribir una estructura ideal a explicar cómo deciden y actúan sus integrantes \citep{larrarteSinFecha,bright2019}. \\

\teoria{Teoría general de sistemas} &
1951--década de 1960 &
Ludwig von Bertalanffy; Daniel Katz; Robert L. Kahn; Fremont Kast; James Rosenzweig &
\aporte{Concibió la organización como sistema abierto compuesto por partes interdependientes.}
\aporte{Articuló entradas, transformación, salidas y retroalimentación, además de totalidad, sinergia, equifinalidad y adaptación.}
\aporte{Obligó a analizar simultáneamente subsistemas técnicos, estructurales, humanos y gerenciales, así como sus intercambios con el entorno.}
Este enfoque superó la lectura aislada de tareas o departamentos \citep{larrarteSinFecha,kastRosenzweig1972,sanchezBanuelos2017}. \\

\rowcolor{itescaSilver!28}
\teoria{Administración estratégica} &
1955--década de 1990 &
Peter Drucker; Alfred Chandler; Igor Ansoff; Henry Mintzberg; Michael Porter &
\aporte{Relacionó objetivos de largo plazo, asignación de recursos, estructura y posición competitiva.}
\aporte{Chandler sostuvo que la estructura sigue a la estrategia; Mintzberg distinguió estrategias deliberadas y emergentes.}
\aporte{Porter aportó marcos para analizar industria, actividades y ventaja competitiva.}
La gestión dejó de concentrarse solo en el funcionamiento interno y pasó a construir cursos de acción frente al mercado y al entorno \citep{sanchezBanuelos2017}. \\

\teoria{Enfoque de contingencia o situacional} &
Finales de 1950--1972; consolidación en las décadas de 1960 y 1970 &
Joan Woodward; Tom Burns; G. M. Stalker; Paul Lawrence; Jay Lorsch; James D. Thompson &
\aporte{Negó la existencia de una única forma óptima de administrar.}
\aporte{Demostró que estructura y prácticas deben ajustarse a tecnología, tamaño, estrategia, incertidumbre y ambiente.}
\aporte{Diferenció configuraciones mecanicistas, apropiadas para estabilidad, y orgánicas, más aptas para cambio e innovación.}
Su pregunta central dejó de ser ``¿cuál es el mejor método?'' para convertirse en ``¿qué configuración es congruente con esta situación?'' \citep{larrarteSinFecha,bright2019}. \\

\rowcolor{itescaSilver!28}
\teoria{Recursos, capacidades y conocimiento organizacional} &
1984--actualidad &
Edith Penrose; Birger Wernerfelt; Jay Barney; Robert Grant; Ikujiro Nonaka; Hirotaka Takeuchi &
\aporte{Explicó la ventaja competitiva desde recursos y capacidades valiosos, escasos y difíciles de imitar o sustituir.}
\aporte{Reconoció el conocimiento como activo intangible que debe crearse, compartirse e incorporarse a procesos, productos y servicios.}
\aporte{Amplió la unidad de análisis desde la empresa aislada hacia aprendizaje, innovación, redes y relaciones interorganizacionales.}
Este enfoque integra adaptación externa con desarrollo de capacidades internas, sin invalidar los aportes previos \citep{sanchezBanuelos2017}. \\

\midrule
\multicolumn{4}{L{24.3cm}}{\scriptsize\textit{Nota.} Elaboración propia a partir de \citet{larrarteSinFecha}, \citet{bright2019} y \citet{sanchezBanuelos2017}. Los periodos indican obras o momentos de consolidación y pueden superponerse; las teorías continúan empleándose de forma combinada según el problema administrativo.} \\
\bottomrule
\end{longtable}
\end{landscape}
\horizontalpagesoff
\restoregeometry
\pagestyle{plain}

\subsection{Lectura evolutiva}

La secuencia permite reconocer cuatro desplazamientos. El primero fue de la improvisación hacia la racionalización del trabajo: Taylor midió tareas, Weber formalizó autoridad y Fayol ordenó funciones directivas. El segundo desplazó la mirada hacia la dimensión humana: las relaciones sociales, la motivación y la racionalidad limitada mostraron que una organización no funciona como una máquina. El tercero abrió sus fronteras: sistemas y contingencia explicaron que los componentes internos son interdependientes y que su eficacia depende del ambiente. El cuarto incorporó estrategia, recursos intangibles y conocimiento para responder a competencia, innovación y cambio acelerado \citep{larrarteSinFecha,sanchezBanuelos2017}.

A juicio del autor, la principal enseñanza no consiste en elegir una teoría como vencedora, sino en diagnosticar qué problema ayuda a comprender cada una. Estandarizar puede mejorar un proceso repetitivo; las reglas pueden proteger imparcialidad; la participación puede elevar compromiso; y una estructura flexible puede favorecer innovación. Sin embargo, aplicar cualquiera de estas soluciones sin considerar personas, tecnología, objetivos y ambiente produce efectos no previstos. La gestión actual requiere, por tanto, integrar eficiencia, legitimidad, bienestar, aprendizaje y capacidad de adaptación.

\clearpage
\section{Conclusiones}

La evolución de la teoría administrativa amplió progresivamente la unidad de análisis: de la tarea al conjunto de la organización, de la estructura a la persona y del interior al entorno. Las escuelas no forman compartimentos excluyentes; representan respuestas históricas que todavía proporcionan herramientas para diagnosticar y decidir. Su valor depende de reconocer tanto su aporte como sus límites.

La tabla muestra que administrar en el presente exige una perspectiva contingente e integradora. La eficiencia operativa carece de sostenibilidad si ignora a las personas; la motivación no basta sin estructura y recursos; y una estrategia pierde pertinencia si no aprende del entorno. Mi postura es que comprender esta trayectoria permite seleccionar prácticas con fundamento, evitar recetas universales y diseñar organizaciones capaces de coordinar, aprender y adaptarse responsablemente. Esta integración es necesaria porque ninguna teoría aislada explica por completo la complejidad organizacional y, en consecuencia, cada decisión administrativa debe ajustarse al problema y al contexto.\footnote{Se utilizó inteligencia artificial como apoyo para organizar ideas, revisar la redacción y editar el documento; no sustituyó el análisis propio. El autor verificó las fuentes y conserva la responsabilidad sobre la interpretación y la versión final.}

\clearpage
\bibliography{fundamentos-de-gestion-administrativa}

\end{document}
